\section{引言}
移动边缘计算（mobile edge computing，MEC）通过将计算服务部署在用户附近，减少任务处理对远端云平台的依赖。无人机（unmanned aerial vehicle，UAV）能够调整服务位置，为临时业务区域和移动用户群提供计算接入。然而，节点可移动并不意味着资源能够随负载即时重分配。当热点在不同区域间迁移时，部分无人机可能持续接收超过本地处理能力的任务，其他无人机则存在空闲算力。仅改善用户与无人机之间的距离，不能消除这种计算负载失衡。

机间协作允许源无人机将部分已上传任务转交其他无人机处理。Xu等研究了多无人机协同计算中的轨迹规划和资源分配问题，说明飞行位置与通信、计算资源之间需要联合考虑\cite{xu2024}。Hao等进一步研究了带任务优先级的协同边缘计算，联合处理任务卸载、资源分配和轨迹设计\cite{hao2024}。上述研究为协作计算提供了重要基础，但不同工作的到达过程、完成判据、资源共享方式和决策时序并不相同。评价一种调度策略时，需要明确这些假设，不能仅凭算法名称或平均时延数字进行比较。

本研究关注连续到达、允许跨帧排队的任务服务过程。任务完整上传后按比例拆分，远端分片须经过机间传输再进入计算队列；只有同一父任务的所有非零分片完成，才能判定该任务完成。在这种处理链中，较短的某一段传输时间不必然带来较低的端到端时延。例如，将一项任务分配给多台无人机可能缓解源节点积压，却也可能使其完成时间受较慢分片限制。与此同时，机间总带宽由当前非空链路共享，新增转发链路会改变其他链路的可用带宽。因此，路由不仅决定计算去向，还改变后续通信竞争。

部署与路由具有不同的调整频率。无人机重新部署需要真实飞行时间和推进能量，不宜随每次任务到达立即移动；任务路由则应及时响应队列和链路变化。多时间尺度协同能够按控制对象区分决策频率。夏玮玮等在无蜂窝无线接入网切片资源分配中采用了多时间尺度协同思路\cite{xia2025}。本研究针对无人机部署期间服务暂停、任务仍然到达的特点，建立共享物理环境下的两层决策过程，而不将两个控制器视为相互独立的环境。

本文采用近端策略优化（proximal policy optimization，PPO）学习大时间尺度部署策略\cite{ppo2017}，采用软演员评论家（soft actor-critic，SAC）学习小时间尺度路由策略\cite{sac2018}。上层根据空间分布和任务压力选择部署方向，下层根据部署后的实际链路和队列选择任务去向。两层均使用由任务驻留量和能耗构成的物理代价，按各自决策间隔聚合反馈。对于周期末服务帧到下一服务决策之间的部署空档，保留真实的等待代价和时间折扣，避免将服务暂停等同于环境终止。

本文的研究内容包括三个方面。建立任务级多无人机协同计算模型，明确上传、一次拆分、机间转发、优先计算与父任务完成之间的关系。建立不同观测接口下的大、小时间尺度决策过程，给出共同目标、动作约束和跨部署边界转移。给出PPO与Direct SAC协同求解过程，区分策略产生的连续分数与实际执行路由，分析动作可行性、价值函数表征和计算复杂度。性能方面需要通过独立训练、配对场景和模块消融检验，不能由上述建模性质直接推出优于其他方法的结论。
